% A11 independent mathematical audit; audit-only; 2026-10-04.
% Audited P11 v1.0 SHA-256:
% 4be8e9b1e64c022b635f3cca9f7f49ffbc5358a2250dda1715cb60e82d9b073b
% All six layer verdicts and OVERALL VERDICT: VERIFIED.
% This source does not modify P11's DRAFT lifecycle status.
\documentclass[11pt]{article}
\usepackage[T1]{fontenc}
\usepackage[utf8]{inputenc}
\usepackage{amsmath,amssymb,amsthm}
\usepackage[margin=1in]{geometry}
\usepackage{array,booktabs,longtable,enumitem}
\usepackage[hidelinks]{hyperref}
\setlength{\emergencystretch}{3em}
\newcommand{\ii}{\mathrm i}
\newcommand{\R}{\mathbb R}
\newcommand{\C}{\mathbb C}
\newcommand{\one}{\mathbf 1}
\newcommand{\lin}[1]{P11 lines #1}
\DeclareMathOperator{\spec}{spec}
\DeclareMathOperator{\diag}{diag}
\DeclareMathOperator*{\argmax}{arg\,max}
\title{A11: Independent audit of\\Hopf spectral conditions and spectral simplicity}
\author{Directed Cyclic Networks research project}
\date{4 October 2026}
\begin{document}
\maketitle
\section{Verdicts and exact-version disposition}
\begin{center}
\begin{tabular}{p{0.73\linewidth}l}
\toprule
GENERAL FREQUENCY / STRICT-CANDIDATE VERDICT & VERIFIED\\
SIMPLICITY / EIGENVECTOR / RESOLVENT VERDICT & VERIFIED\\
CANDIDATE-SURFACE VERDICT & VERIFIED\\
REALIZABILITY-CRITERION VERDICT & VERIFIED\\
FIVE-NODE CANDIDATE VERDICT & VERIFIED\\
SCOPE / P12-P13 HANDOFF VERDICT & VERIFIED\\
\midrule OVERALL VERDICT & VERIFIED\\
\bottomrule
\end{tabular}
\end{center}
\noindent\textbf{Required corrections: none.}
The six verdicts concern the precise spectral claims in P11 v1.0. They do
not certify a parameter crossing, a Hopf bifurcation, a Lyapunov coefficient,
a periodic orbit, a rotating wave, or a first loss of stability.

\medskip\noindent{\raggedright\textbf{PROPOSED SCOPE CLARIFICATION RECOMMENDATION:}\par}
\noindent\textbf{ACCEPT AS PLANNING CLARIFICATION.}\par\smallskip
This recommendation is separate from the mathematical verdict. No frozen
file is amended or silently superseded here.

\paragraph{Exact audited source.}
\path{proofs/P11_hopf_spectral_conditions.tex}, P11 v1.0,
3 October 2026, 25,914 bytes and 600 lines. Its full-file SHA-256 is
\begin{center}\ttfamily\small
4be8e9b1e64c022b635f3cca9f7f49ffbc53\\
58a2250dda1715cb60e82d9b073b
\end{center}
The two displayed segments concatenate to the full hash. The manifest's
last P11 CURRENT entry designates this exact file and hash. Earlier P11
placeholders in other snapshots are not audit inputs. All original files
remain unchanged; the audit verdict does not promote the proof's lifecycle.

\paragraph{Method.}
This is an adversarial analytical reconstruction from the canonical field,
current P02/P04/P05, and conditional P07 identities. Every source line is
covered by the range ledger below. Exact and finite numerical checks are
separate corroboration. The developer validation package was inventoried
and hashed but its checker was neither read as a derivation nor executed.
No literature search or later-module mathematics was undertaken.

\section{Current-file and dependency authentication}
\label{sec:auth}
The working root is \path{cyclic_network_project/}. Both copies of
\path{cyclic_network_project_P11_draft.zip} have SHA-256
\begin{center}\ttfamily\small
8513c67effa28b825b52052e2c9333282f66\\
3b1031610ea150e58aeb16b58ad5.
\end{center}
All 200 file members pass CRC and match the working files byte-for-byte.
The P10 audited baseline has 199 members and SHA-256
\begin{center}\ttfamily\small
e0ecae17063fced8ed0a6ebb6fcf44456d5\\
c22065f51059274a0231659dcaaf5.
\end{center}
Exactly five baseline members differ: P11, manifest, registry, compendium
source and compendium PDF. The proposed clarification is the sole added
member. All requested controls and prior proof/audit sources are present.
The 31-item input inventory records names, full hashes, byte and line
counts; proof mathematical-body hashes are also recorded. The 200-item
baseline inventory permits a final preservation check independently of
any lifecycle prose.

The authenticated controls are the manifest, claim registry, frozen theorem
package, final novelty verdict, novelty decisions, open decisions,
dependency map, proofs compendium, shared macros, canonical model, and
notation/assumptions. Current proofs and audits are P01/A01, P02/A02,
P03/A03, P04/A04, P05/A05, P07/A07, P08/A08, P09/A09 and P10/A10.
Each current proof's computed full hash appears in its manifest closure.
P01/P02 and P05/P08 retain their recorded E01 closures; P03/P04/P07/P09/P10
retain metadata-only closures and their body hashes. Original audits keep
their historical verdicts. A02's source exists: its old missing-source
header is superseded by the manifest's restoration record. A04/A05's
historical validation-provenance qualifications are not erased.

\paragraph{Actual premises, not merely header declarations.}
The dependency map directly lists P02/P04/P05 $\to$ P11, and separately
requires P07 for five-node use. The body respects this graph:
\begin{itemize}[nosep]
\item P02 supplies the nonzero equilibrium equation, scalar polynomial,
root classification and scalar bracket (lines 239--248, 475--490,
536--544).
\item P04 supplies the oriented Fourier spectrum, conjugacy, normality,
label-count multiplicity and uniform-mode identity (41--50, 141--166,
193--207, 245--260, 451--452, 537--539).
\item P05 supplies complete selection, reachable representatives, strict
cell endpoints and the concave equilibrium chart (54--76, 141--147,
291--300, 322--387).
\item P07 is used only in the five-node subsection (391--398, 453--470).
Its conditional role is preserved.
\end{itemize}
The inherited P03 coefficient formula is restated through P04, not used
as a new direct DAG edge. P08/P09/P10 are authenticated context only.
The two-positive-root argument actually factors P02's polynomial, so it
does not conceal a P10 premise. P12/P13 are future recipients, and
P14--P18 supply no results. The P13 placeholder is inspected solely as a
convention requirement, not as a mathematical input. Elementary linear
algebra and the local implicit/regular-level-set theorems supply the
standard tools explicitly used below.

\section{Layer A: general frequency and strict selection}
\label{sec:A}
The original predecessor convention is $(Pv)_i=v_{i-1}$ and
$v_i^{(m)}=n^{-1/2}e^{2\pi\ii mi/n}$. Thus the predecessor contributes
$Ce^{-\ii\theta_m}$, the positive successor-two term contributes
$Ge^{2\ii\theta_m}$, and the negative predecessor-two term contributes
$-He^{-2\ii\theta_m}$. The sign of the last imaginary contribution is
positive. Consequently
\[
 r_m=d+C\cos\theta_m+(G-H)\cos2\theta_m,\qquad
 \omega_m=-C\sin\theta_m+(G+H)\sin2\theta_m.
\]
These agree literally with current P04; no $G/H$ phase reversal occurs.

P05 gives $1\le j\le\lfloor(n+1)/4\rfloor<n/2$ for every reachable
nonuniform representative. The strict inequality follows from $n>1$.
For completeness, if $n=4r$, the largest representative is $r<2r$;
if $n=4r+2$, it is $r<2r+1$. These cover every even size. If $n=4r+1$
it is $r$, and if $n=4r+3$ it is $r+1$, again strictly below $n/2$.
Hence $0<\theta_j<\pi$, $\sin\theta_j>0$, and the two labels
$j,n-j$ are distinct. The Nyquist singleton is never reachable.

Using $\sin2\theta=2\sin\theta\cos\theta$ gives independently
\[
 \omega_j=\sin\theta_j\{2(G+H)u_j-C\},\qquad u_j=\cos\theta_j.
\]
Since $\sin\theta_j>0$, zero frequency is equivalent to
$2(G+H)u_j=C$. For $u_j\le0$ the bracket is at most $-C<0$,
so the frequency is strictly negative. For $u_j>0$, the zero locus is
exactly $G+H=C/(2u_j)$, with the claimed sign on either side.
There is no decimal or asymptotic premise.

The definition requires $j\ne0$, the complete selected set
$\mathcal S=\{j\}$, $r_j=0$ and $\omega_j\ne0$. It excludes every
maximal P05 tie, not every accidental equality among stable real parts.
For every other representative $t$, completeness gives $r_t<0$.
Every Fourier label belongs to exactly one representative class, so all
labels outside $L_j=\{j,n-j\}$ are strictly left. This includes label
zero, all conjugate partners, and the even-$n$ singleton. Repeated
numerical eigenvalues among stable labels do not change the conclusion.
P04 conjugacy then gives precisely $\lambda_j=\ii\omega_j$ and
$\lambda_{n-j}=-\ii\omega_j$ on the imaginary axis.

\noindent\textbf{GENERAL FREQUENCY / STRICT-CANDIDATE VERDICT: VERIFIED.}

\section{Layer B: simplicity, vectors, real plane and inverse gates}
\label{sec:B}
If a label outside $L_j$ produced $\lambda_j$, its real part would be
zero, contrary to strict selection. The other label in $L_j$ gives the
opposite value and is distinct because $\omega_j\ne0$. Exactly one
label produces $\lambda_j$. Interchanging the two labels gives the
same result for $\lambda_{n-j}$. P04's multiplicity theorem makes each
value algebraically and geometrically simple over $\C$. There are two
simple eigenvalues, not one eigenvalue of multiplicity two.

There are $n-2\ge3$ other labels. This finite nonempty set has a
strictly negative maximum real part, so
\[
 \eta=-\max_{\ell\notin L_j}r_\ell>0,
 \qquad \Re\lambda_\ell\le-\eta\quad(\ell\notin L_j).
\]
The critical Fourier span is conjugation-invariant and has a real form of
dimension two. The complementary Fourier span is likewise
conjugation-invariant, with real dimension $n-2$, and is stable. It may
contain repeated stable eigenvalues; diagonalizability remains valid.

Both $0$ and $2\ii\omega_j$ have zero real part, excluding equality
with every stable eigenvalue. Zero differs from both critical values.
Equality $2\ii\omega_j=\ii\omega_j$ would require $\omega_j=0$;
equality $2\ii\omega_j=-\ii\omega_j$ would require $3\omega_j=0$.
Thus both $J^{-1}$ and $(2\ii\omega_jI-J)^{-1}$ exist, for either
sign of $\omega_j$. No nonlinear formula enters this argument.

\paragraph{Adjoint and basis audit.}
Write $J=Q\Lambda Q^*$ with unitary $Q$. Then
$J^*=Q\overline\Lambda Q^*$; for $v_H=v^{(j)}$ one has
\[
 Jv_H=\ii\omega_jv_H,\qquad
 J^*v_H=-\ii\omega_jv_H,\qquad v_H^*v_H=1.
\]
Thus $p_H=v_H$ is correct for the actual pairing $\langle p,v\rangle=p^*v$.
This also agrees with the current B07 resolution of the opposite-sign
adjoint equation. The TODO P13 source requires reconciliation of this
pairing and Taylor factors but states no competing adopted formula.
There is no convention mismatch; P13 still owns its Taylor, factorial,
time and coefficient normalizations.

Put $v_H=A+\ii B$. Orthogonality of $v_H$ and $\overline v_H$
yields $v_H^{\mathsf T}v_H=0$, while $v_H^*v_H=1$. Taking real
and imaginary parts gives
\[
 \|A\|^2-\|B\|^2=0,\quad A^{\mathsf T}B=0,\quad
 \|A\|^2+\|B\|^2=1.
\]
Therefore $e_c=\sqrt2 A,e_s=\sqrt2 B$ are orthonormal. Real and
imaginary parts of the eigenvector equation give
\[
 Je_c=-\omega_j e_s,\quad Je_s=\omega_j e_c,
 \qquad [J]_{(e_c,e_s)}=
 \begin{pmatrix}0&\omega_j\\-\omega_j&0\end{pmatrix}.
\]
The columns have exactly the stated orientation. For a nonzero coefficient
vector $(z_c,z_s)$ its angle has derivative
$(z_c\dot z_s-z_s\dot z_c)/(z_c^2+z_s^2)=-\omega_j$.
This is a linear coefficient-plane statement, not nonlinear propagation.

\paragraph{Excluded degeneracies.}
Two different representative classes at real part zero supply additional
critical labels, whether or not their numerical eigenvalues coincide.
The single-simple-pair hypothesis fails. P11 gives no named
multiple-criticality bifurcation classification. If a genuine pair has
$r_j=\omega_j=0$, both labels produce zero. P04 makes this eigenvalue
semisimple, with algebraic and geometric multiplicity equal to the total
number of zero labels. Under strict representative selection that number
is exactly two and the real kernel dimension is two. Without strictness,
all other zero labels must be counted. Neither case is admitted as a
strict spectral Hopf candidate.

\noindent\textbf{SIMPLICITY / EIGENVECTOR / RESOLVENT VERDICT: VERIFIED.}

\section{Layer C: regular candidate surfaces}
\label{sec:C}
P02's nonzero equilibrium relation gives precisely P05's open chart
\[
 a=(k+h)x-\frac{c-g}{1+x}>0,\qquad x,k,c,g,h>0.
\]
Substitution into the unsimplified diagonal coefficient gives
\begin{align*}
 d&=a-2kx-\frac{cx}{(1+x)^2}-\frac{g}{1+x}-hx\\
  &=-kx-\frac{c}{1+x}-\frac{cx}{(1+x)^2}
   =-kx-\frac{c(1+2x)}{(1+x)^2}.
\end{align*}
Let $v_j=2u_j^2-1=\cos2\theta_j$, a scalar distinct from $v^{(j)}$.
It follows that
\[
 R_j=-kx-\frac{c(1+2x)}{(1+x)^2}+\frac{cu_j}{1+x}
          +x\left(\frac{g}{(1+x)^2}-h\right)v_j=r_j.
\]
At fixed $x,c,g,h$, $\partial_kR_j=-x<0$, so $R_j=0$ is the
regular graph
\[
 k=k_H=\frac1x\left[-\frac{c(1+2x)}{(1+x)^2}
       +\frac{cu_j}{1+x}
       +x\left(\frac{g}{(1+x)^2}-h\right)v_j\right].
\]
Physical admissibility still requires $k_H>0$ and
$a_H=(k_H+h)x-(c-g)/(1+x)>0$. These restrictions cannot be dropped.
Inside a strict selected cell, nonzero frequency is an additional open
restriction on this regular level set. The claim is not that every
reachable cell contains a nonempty portion; Layer D addresses exactly that.

P05 proves that $G-H\ge0$ permits only the uniform maximizer. Therefore
any nonuniform selected representative forces $D_c=H-G>0$ and
$h-g/(1+x)^2=D_c/x>0$. For the map
$\Phi(x,k,c,g,h)=(a,k,c,g,h)$ the determinant is
\[
 \det D\Phi=k+h+\frac{c-g}{(1+x)^2}
       =k+\frac{D_c}{x}+\frac{c}{(1+x)^2}>0.
\]
The local inverse function theorem transfers the regular surface to
physical parameter space along the specified branch. This does not
identify different branches globally and does not prove global injectivity.
The derivative in $k$ above lets $a$ follow the chart; it is not a
crossing derivative along the later fixed-$a$ path.

For $u_j>0$,
\[
 \partial_k\omega_j=0,\qquad
 \partial_g\omega_j=\frac{2xu_j\sin\theta_j}{(1+x)^2}>0.
\]
The determinant of the $(k,g)$ minor of $D(R_j,\omega_j)$ is
\[
 -\frac{2x^2u_j\sin\theta_j}{(1+x)^2}\ne0.
\]
Any physically admissible nonempty joint level set is locally codimension
two in branch coordinates and codimension one in $R_j=0$. Physical
transfer is conditional on the same concave local chart. For $u_j\le0$
the joint set is empty by the strict sign already proved. A statement of
rank does not itself assert nonemptiness or a bifurcation unfolding.

\noindent\textbf{CANDIDATE-SURFACE VERDICT: VERIFIED.}

\section{Layer D: independent reconstruction of realizability}
\label{sec:D}
Fix a reachable nonuniform representative and write
\[
 w_j=1-2u_j^2,\quad
 L=\max(0,\beta_j),\quad U=\beta_{j-1},\quad
 B=\frac{w_j}{4(1-u_j)},\quad z=\frac{C}{4D_c}.
\]
P05's cosine grid is strictly decreasing. Its positive strict cell is
$L<z<U$. The endpoint $\beta_j$ exists because
$j\le j_{\max}<N=\lfloor n/2\rfloor$. Reachability gives $U>0$;
strict midpoint decrease gives $L<U$. In particular a zero lower midpoint
is excluded by $z>0$, not included as an attainable tie.

\paragraph{Necessity from positive physical coefficients.}
Use $c=C(1+x)$ and $G-H=-D_c$ in the independently derived $R_j$.
Then $R_j=0$ is equivalent to
\[
 kx=D_cw_j-C\left(1-u_j+\frac{x}{1+x}\right),\qquad
 \frac{kx}{C}=\frac{w_j}{4z}-1+u_j-\frac{x}{1+x}.
\]
Since $k,x,C>0$ and $u_j<1$, this forces
\[
 w_j>0,\qquad \frac{w_j}{4z}>1-u_j+\frac{x}{1+x}.
\]
The first condition is exactly $|u_j|<1/\sqrt2$. Equality $w_j=0$
cannot work with $k>0$. Dropping the positive term $x/(1+x)$ yields
$z<B$. Combining this with $z>L$ gives $L<B$. Thus every physical
strict spectral candidate must satisfy both claimed inequalities.

\paragraph{Sufficiency, including all physical restrictions.}
Assume $w_j>0$ and $L<B$. Both $U$ and $B$ exceed $L$, so choose
\[
 L<z<\min(U,B).
\]
Then $\delta=w_j/(4z)-(1-u_j)>0$. For example,
$x=\delta/[2(1+\delta)]>0$ satisfies $x/(1+x)<\delta$;
this makes the fixed-vertex inequality strict. More generally, for any
prescribed $x,z$ satisfying it, choose any $C>0$, put $D_c=C/(4z)$,
and set
\[
 k=\frac{D_cw_j-C(1-u_j+x/(1+x))}{x}>0.
\]
For a free $G>0$, define
\begin{align*}
 c&=C(1+x),&g&=\frac{G(1+x)^2}{x},&h&=\frac{G+D_c}{x},\\
 a&=kx+D_c-C+G\left(1+\frac{1+x}{x}\right).
\end{align*}
The first four coefficients are positive. Their effective coefficients
are exactly $C,G,H=G+D_c$, hence the selected vertex is still $z$.
Directly,
\[
 (k+h)x-\frac{c-g}{1+x}
 =kx+G+D_c-C+\frac{G(1+x)}x=a.
\]
This proves the equilibrium equation, and the definition of $k$ proves
$R_j=0$. The coefficient of $G$ in $a$ is positive, so sufficiently
large $G$ gives $a>0$ without moving $k,z,x$ or $R_j$.

The frequency bracket is
\[
 2(G+H)u_j-C=2(D_c+2G)u_j-C.
\]
For $u_j\le0$ it is strictly negative. For $u_j>0$ its only possible
zero is
\[
 G_0=\frac{C}{4u_j}-\frac{D_c}{2}.
\]
At most one positive value is forbidden; arbitrarily large positive $G$
with $a>0$ can avoid it. This proves sufficiency both for the
fixed-$x,z$ statement and for the full-cell statement. All parameters
remain finite; no limit point at $x=0$ is used as a physical realization.
The necessary-and-sufficient criterion is therefore exactly
\[
 \boxed{w_j>0\quad\hbox{and}\quad\max(0,\beta_j)<
                   \frac{w_j}{4(1-u_j)}.}
\]

\paragraph{Exact obstruction checks.}
For $n=8,j=1$, $u_1=\cos(\pi/4)=1/\sqrt2$ and $w_1=0$.
The equality for $kx$ would have a negative right side. The cell can
realize dominance, but no physical zero spectral abscissa in that cell.
For $n=12,j=2$, $u_2=1/2$, $u_3=0$, and
\[
 w_2=\frac12,\qquad \beta_2=\frac{u_2+u_3}{2}=\frac14,
 \qquad L=\frac14=B.
\]
The upper endpoint is $\beta_1=(\sqrt3+1)/4>1/4$; the cell is
nonempty. Every strictly interior vertex has $z>1/4$, contrary to the
necessary $z<B=1/4$. Even the lower boundary would require
$x/(1+x)<0$ for positive $k$. These are genuine obstructions, not a
midpoint indexing error. Neither contradicts P05's dominance realization.
P11 replaces, rather than retains, the false optional universal assertion.

\noindent\textbf{REALIZABILITY-CRITERION VERDICT: VERIFIED.}

\section{Layer E: exact five-node candidate and spectral handoff}
\label{sec:E}
Set $\phi=(1+\sqrt5)/2$ and
$\sigma=\sqrt{10-2\sqrt5}/4=\sin(\pi/5)>0$. Then
$\phi^2=\phi+1$, $\phi^{-1}=\phi-1$, $1<\phi<2$.
For
\[
 x_H=1,\quad c=2,\quad g_H=16,\quad h=6,\quad
 k_H=\frac{3\sqrt5-5}{4}=\frac{3\phi-4}{2},\quad
 a_H=\frac{47+3\sqrt5}{4}=\frac{3\phi+22}{2},
\]
positivity of $k_H$ follows from $5/3<\sqrt5$, because
$25/9<5$. Positivity of the other four coefficients is immediate.
In the original canonical equation, before polynomial reduction,
\[
 F_i(\one)=a_H-k_H+\frac22-\frac{16}{2}-6
          =a_H-k_H-13=0.
\]
This verifies all five components of the actual DM field.

The effective values are $s=2,C=1,G=4,H=6,T=-2,S=10$, with
\[
 d=a_H-2k_H-\frac12-8-6
   =-k_H-\frac32=\frac{1-3\phi}{2}.
\]
P07's exact phase data are $u_1=(\phi-1)/2$, $u_2=-\phi/2$,
$\sin\theta_1=\phi\sigma$, $\sin\theta_2=\sigma$,
$\sin2\theta_1=\sigma$, $\sin2\theta_2=-\phi\sigma$.
Substitution, without using only an interval table, gives
\begin{align*}
 r_1&=\frac{1-3\phi}{2}+\frac{\phi-1}{2}+\phi=0,\\
 r_0&=\frac{1-3\phi}{2}+1-2=-\frac{3\phi+1}{2},\\
 r_2&=\frac{1-3\phi}{2}-\frac\phi2-(\phi-1)
       =-\frac{3\sqrt5}{2}.
\end{align*}
Their difference $r_0-r_2=(3\phi-4)/2=k_H>0$ proves
$r_1>r_0>r_2$ exactly. The imaginary parts are
\[
 \omega_H=-\phi\sigma+10\sigma=\sigma(10-\phi)>0,
 \qquad\omega_2=-\sigma-10\phi\sigma.
\]
The exact inequality $10-\phi>8$ certifies nonzero frequency. Thus
\begin{align*}
 \lambda_0&=-\eta_H,&\eta_H&=\frac{3\phi+1}{2},\\
 \lambda_1&=\ii\omega_H,&\lambda_4&=-\ii\omega_H,\\
 \lambda_2&=-\frac{3\sqrt5}{2}-\ii\sigma(1+10\phi),
 &\lambda_3&=\overline{\lambda_2}.
\end{align*}
The general strict criterion proves individual simplicity of labels $1,4$
and excludes all other imaginary-axis labels. The gap is exactly
$-\max(r_0,r_2)=\eta_H$, not merely a lower estimate.

\paragraph{Upper-root identity from P02 only.}
Put $A=k_H+h=(3\phi+8)/2$, $b=a_H-A=7$ and
$\mu=a_H+2-16=(3\phi-6)/2<0$. Since $A-b=\mu$,
\[
 Q(X)=AX^2-bX-\mu=A(X-1)(X-x_-),\qquad
 x_-=-\frac\mu A=\frac{3(2-\phi)}{3\phi+8}.
\]
The numerator and denominator are positive, and their difference is
$6\phi+2>0$, proving $0<x_-<1$. The discriminant equals
$A^2(1-x_-)^2>0$. Thus there are two distinct positive roots and $1$
is precisely the upper root. Merely knowing $b>0,\mu<0$ would not
have established this without checking the discriminant; P11 checks it.

\paragraph{Local branch, without a path derivative.}
Freeze $\pi(g)=(a_H,k_H,2,g,6)$, with $g_H=16$ and $j=1$.
The scalar bracket is
\[
 B(X,g)=a_H-(k_H+6)X+\frac{2-g}{1+X}.
\]
At $(1,16)$, $B=0$ and
\[
 B_X=-A-\frac{2-16}{4}=-A+\frac72
      =-\frac{3\phi+1}{2}=-\eta_H<0.
\]
P04's uniform eigenvalue is $f_X=B+XB_X$; the equality
$B_X=\lambda_0$ here uses both $B=0$ and the particular state $X=1$.
It is not an incorrectly generalized identity for arbitrary $X$.
The scalar implicit-function theorem gives a unique smooth local root
near $1$, with positive state and parameters. Distinctness, positivity and
ordering of the two roots persist by continuity, so this is $x_+(g)$.
No $dx_+/dg$ or eigenvalue derivative along this path is computed.

\paragraph{Exact vectors and resolvents.}
The required vector is the positive-exponent Fourier column
\[
 v_H=\frac1{\sqrt5}(1,e^{2\pi\ii/5},e^{4\pi\ii/5},
 e^{6\pi\ii/5},e^{8\pi\ii/5})^{\mathsf T},\qquad p_H=v_H.
\]
With node order $0,\ldots,4$, Layer B verifies
$J_Hv_H=\ii\omega_Hv_H$, $J_H^*p_H=-\ii\omega_Hp_H$ and
$p_H^*v_H=1$. Its real basis has precisely the stated block and
coefficient-plane angular velocity $-\omega_H$.
The complete orthogonal spectral projectors give
\begin{align*}
 J_H^{-1}&=\sum_{m=0}^4\lambda_m^{-1}v^{(m)}v^{(m)*},\\
 (2\ii\omega_H I-J_H)^{-1}
 &=\sum_{m=0}^4(2\ii\omega_H-\lambda_m)^{-1}v^{(m)}v^{(m)*}.
\end{align*}
All denominators are nonzero by Layer B. In particular,
\[
 \lambda_0=-\eta_H\ne0,\qquad
 2\ii\omega_H-\lambda_2
 =\frac{3\sqrt5}{2}+\ii\sigma(21+8\phi)\ne0,
\]
because $2(10-\phi)+(1+10\phi)=21+8\phi$. This supplies spectral
data only; no nonlinear interaction channel is evaluated.

\noindent\textbf{FIVE-NODE CANDIDATE VERDICT: VERIFIED.}

\section{Layer F: literal ledger, scope and planning recommendation}
\label{sec:F}
\begin{longtable}{p{0.08\linewidth}p{0.84\linewidth}}
\toprule Item&Independent disposition\\\midrule\endhead
H1&Established: all five physical coefficients are strictly positive.\\
H2&Established: each original-field component vanishes exactly at $\one$.\\
H3&Established: the fixed $g$-path has a unique local positive upper root near $1$.\\
H4&Established: labels $1,4$ give the positive-frequency pair $\pm\ii\omega_H$.\\
H5&Established: both numerical critical eigenvalues are separately simple.\\
H6&Established: every other eigenvalue has real part at most $-\eta_H<0$.\\
H7&Established: zero is outside the spectrum.\\
H8&Established: $2\ii\omega_H$ is outside the spectrum.\\
H9&Established: exact normalized right/adjoint vectors and real-plane orientation.\\
H10&Correctly unestablished: transversality; no state or eigenvalue derivative along $g$.\\
H11&Correctly unestablished: first Lyapunov coefficient and nonlinear nondegeneracy.\\
H12&Correctly unestablished: periodic orbit, branch stability and rotating wave.\\
\bottomrule
\end{longtable}
The branch-coordinate $k$ derivative and the partial frequency derivative
in the surface proof do not assert H10. The scalar $B_X$ derivative proves
regularity of an equilibrium branch; it is not $dx_+/dg$. No second or
third field derivative, normal-form tensor, Fourier nonlinear interaction,
Lyapunov value/sign, or periodic-branch stability appears in P11.
The word Hopf is consistently qualified as spectral candidate or future
package. No earlier paragraph overrides H10--H12.

\paragraph{Separate clarification review.}
The proposed note correctly identifies frozen attempt-order row 10,
``certify first crossing and exclusions,'' and novelty-decisions S11,
which seeks an open first-loss regime along a specified path. The latter
also requires a nonempty regime, a regular positive branch, complete
spectral exclusions, and later transversality/nonlinear nondegeneracy.
An isolated exact point and a regular spectral surface do not by themselves
complete that target.

The proposed clarification allocates P11 the spectral candidate, its
regular branch and exclusion data; P12 the moving-equilibrium crossing
derivative; and P13 the nonlinear nondegeneracy and assembled Hopf theorem.
A first-loss assertion still requires an explicit interval/path statement.
This is consistent with the actual dependency map, which already assigns
spectral conditions to P11, transversality to P12, and the assembled theorem
to P13. The exact five-node point and Layer C prove the asserted nonempty
regular candidate surface locally. The cell criterion refines an optional
existence question and does not retract any P05 theorem.

Adoption would clarify the allocation and wording of planning obligations;
it would not change direct or conditional dependency edges, canonical
conventions, novelty ratings, the assembled first-Hopf target, or currently
permitted manuscript claims. The stronger first-loss/open-regime target
remains unproved and is not discharged by this audit. Exact-model priority
remains unclosed under the existing bounded novelty record. No new novelty
assessment or literature survey is made. The clarification should remain
an explicit planning action at closure, rather than being treated as a
mathematical theorem or an automatic frozen-file edit.

\noindent\textbf{Recommendation: ACCEPT AS PLANNING CLARIFICATION.}
\par\noindent\textbf{SCOPE / P12-P13 HANDOFF VERDICT: VERIFIED.}

\section{Adversarial checks and validation provenance}
\label{sec:checks}
The attempted failure modes were treated analytically, not merely by
searching for numerical counterexamples:
\begin{enumerate}[leftmargin=2em]
\item Pair reachability: all four residue classes of $n$ and the Nyquist
endpoint were checked in Layer A.
\item Extra imaginary labels: complete representative selection covers all
labels, including uniform and even singleton labels.
\item False simplicity: label equality was separated from numerical-value
equality; the latter is excluded using strict real parts and nonzero frequency.
\item A $2\ii\omega$ pole: stable and both critical values were checked separately.
\item Wrong adjoint sign: the actual unitary diagonalization and Hermitian
pairing give the opposite-sign adjoint equation.
\item Failed physical projection: strict nonuniform selection forces
$D_c>0$, making the chart determinant strictly positive.
\item Missing realizability restrictions: necessity includes $k>0$, and
sufficiency separately enforces $a>0$ and avoids the sole possible frequency zero.
\item The $(8,1)$ obstruction survives the boundary $w=0$ exactly.
\item The $(12,2)$ obstruction uses $\beta_2$, not $\beta_1$; equality is insufficient.
\item The five-node tuple is substituted into the original field first.
\item Every five-node eigenvalue and imaginary sign is reconstructed from
exact phase data, then independently checked in the direct matrix.
\item The upper-root identity is established by exact factorization and
positive discriminant, without P10.
\item The proposed note is checked against the actual stronger planning
language; accepting it does not certify first loss or reduce the assembled target.
\end{enumerate}

\paragraph{An explicit physical frequency-degenerate check.}
This independent exclusion test uses $n=5,j=1$, $x=1/10$, $C=1$,
$D_c=1$, $G=(\phi-1)/2$ and $z=1/4$. The strict cell lower endpoint
is zero and its upper endpoint is $(1+\phi)/4$, so this vertex is strict.
Here $u=(\phi-1)/2$, $w=\phi/2$ and
\[
 kx=\phi-\frac{35}{22}>0,
\]
since $\sqrt5>24/11$ follows by squaring ($605>576$).
Use the physical construction in Layer D. Then $a=kx+12G>0$,
all other coefficients are positive, $R_1=0$, but
$2(D_c+2G)u-C=2\phi u-1=0$. There is an exact two-dimensional
real kernel with stable complement, not a strict spectral Hopf candidate.
This demonstrates a genuine physical excluded point, not a fabricated
free-coefficient example. It also witnesses that the conditional rank-two
statement need not be vacuous.

\paragraph{An explicit critical tie check.}
For $n=20$, the neighboring representatives $4,5$ have cosines
$u_4=(\phi-1)/2$ and $u_5=0$. Choose $C=1$, $D_c=\phi$,
$z=1/(4\phi)=(u_4+u_5)/2$, $x=1/10$ and $G=1$.
Then $kx=\phi-12/11>0$ and the same physical construction gives
$a=2\phi+109/11>0$. Both representative real parts equal zero,
and P05's midpoint rule excludes other representatives. Four Fourier
labels are critical. This fails strict selection regardless of numerical
frequency coincidences; no named bifurcation is inferred.

\paragraph{Independent exact arithmetic.}
The new checker uses rational arithmetic in
$\mathbb Q(\phi,t)$ with $\phi^2=\phi+1$ and
$t^2=(\phi-3)/4$, interpreted as $t=\ii\sigma$.
The fifth root is $z=(\phi-1)/2+\phi t$. Checking unnormalized
columns $w_i^{(m)}=z^{im}$ gives $w^{(m)*}w^{(m)}=5$;
division by $\sqrt5$ yields the project normalization exactly.
The script performs 144 exact assertions, including the original residual,
direct $5\times5$ matrix eigenvectors, Gram matrix, adjoint residuals,
real basis identities, both spectral-inverse products, mode denominators,
branch polynomial and discriminant identities, exact obstruction constants,
and two exact null columns at the excluded point. Exact inequalities are
proved above; floating-point readings are not used to certify them.

\paragraph{Finite corroboration and limits.}
With Python 3.12.14 and NumPy 2.3.5, the separate checker examined all
646 reachable nonuniform cells for $5\le n\le72$: 179 satisfy the
criterion and 467 do not. It constructed 358 strict physical candidates,
ran 9,340 obstruction probes, and constructed 111 physical critical ties.
A separately seeded raw-parameter search made 1,500 trials and found 147
admissible strict candidates, all consistent with necessity. There were
206 direct matrix-spectrum comparisons, including odd and even sizes.
Maximum scaled direct-spectrum discrepancy was $3.366\times10^{-15}$;
right and adjoint residuals were below $3.167\times10^{-15}$, original
field residual below $1.074\times10^{-16}$. The chart/uniform-identity
check had maximum scaled cancellation error $1.417\times10^{-12}$.
The five-node reading is approximately
\[
 -2.9270509831,\quad \pm4.9267960066\ii,\quad
 -3.3541019662\pm10.0983504152\ii.
\]
These finite checks corroborate, but do not prove, all-$n$ claims or certify
intervals. Floating-point cell equality uses a tolerance; the two specified
boundary obstructions are checked exactly. The saved JSON contains the
actual counts and residuals; no historical developer statistics are claimed
as independently reproduced.

\paragraph{Developer validation authentication.}
\path{output/p11/P11_validation_support.zip} has SHA-256
\begin{center}\ttfamily\small
3fc91d1788cbbbec0ecdd0f8df972481a71\\
33e6a87ce0f44bdcacf187d8ad18d.
\end{center}
Its CRC and all 36 members were inventoried; the original remains unchanged.
This establishes package identity only. The checker and historical counts
were not rerun or adopted as proof. The audit's new script and results are
independently authored. Compilation and visual inspection likewise do not
supply mathematical verification.

\section{Complete source-line review}
\label{sec:lines}
The following ranges partition all 600 lines without gaps or overlaps.
Every substantive range receives exactly one permitted disposition. Blank
lines, labels and environment delimiters are attached to their neighboring
substantive range. PASS means valid within the stated scope, not a proof of
an excluded later result. All findings refer only to the authenticated hash.
\begin{longtable}{p{0.10\linewidth}p{0.10\linewidth}p{0.70\linewidth}}
\toprule Lines&Verdict&Finding\\\midrule\endhead
1--18&PASS&Version, inputs and exclusions match the manifest; no numerical premise.\\
19--31&PASS&Spectral scope is explicit; no crossing, orbit or propagation assertion.\\
32--40&PASS&Positive equilibrium and effective coefficients use the canonical domain.\\
41--53&PASS&P04 phases, signed imaginary parts and unitary diagonalization agree exactly.\\
54--64&PASS&Complete representative maximizers include all ties and singleton classes.\\
65--79&PASS&Reachable bound proves a genuine pair; all even sizes exclude Nyquist.\\
80--96&PASS&Frequency factorization and both cosine-sign cases are exact.\\
97--103&PASS&Double-angle identity and positive sine justify every equivalence.\\
104--114&PASS&Strict means one complete representative, zero real part, nonzero frequency.\\
115--125&PASS&Exactly two distinct simple imaginary eigenvalues; every other label is left.\\
126--139&PASS&Finite positive gap, real dimensions and both inverse gates are correctly stated.\\
140--148&PASS&Representative-to-label transfer accounts for uniform and Nyquist modes.\\
149--159&PASS&Numerical simplicity follows from label exclusion and P04 multiplicity.\\
160--168&PASS&Noncritical set is nonempty; both critical-value comparisons are valid.\\
169--182&PASS&Same normalized Fourier column is the opposite-sign Hermitian adjoint.\\
183--191&PASS&Ordered real block and signed coefficient-plane angular velocity are correct.\\
192--198&PASS&Unitary diagonalization verifies the current pairing convention.\\
199--205&PASS&DFT orthogonality yields equal real/imaginary norms and a real basis.\\
206--214&PASS&Column signs are correct; negative signed frequency is handled explicitly.\\
215--223&PASS&Critical ties fail the single-pair hypothesis; real-part ties need not be eigenvalue ties.\\
224--234&PASS&Zero multiplicity is exactly the zero-label count; strict zero pair has real kernel dimension two.\\
235--244&PASS&P02 equilibrium chart is exactly the P05 open physical domain.\\
245--251&PASS&Diagonal simplification cancels the g terms; scalar cosine notation is distinguished.\\
252--261&PASS&R\_j equals the selected modal real part with all physical coefficients.\\
262--272&PASS&Partial k derivative and solved graph have correct signs.\\
273--280&PASS&Regular surface retains k and a positivity, strict selection and nonzero frequency.\\
281--286&PASS&Rank-two assertion is conditional on existence and correct physical transfer.\\
287--301&PASS&P05 forces concavity; determinant is positive; transfer is local only.\\
302--305&PASS&Branch-coordinate derivative is expressly separated from the frozen physical path.\\
306--318&PASS&Frequency partials and negative nonzero two-by-two determinant are exact.\\
319--334&PASS&Cell criterion uses the correct lower midpoint and strict inequality.\\
335--347&PASS&Fixed vertex/state iff claim is correct; both obstruction examples are valid.\\
348--361&PASS&Positive k forces w positive and z below B; all endpoints exist with correct indexing.\\
362--371&PASS&Both upper bounds exceed the lower bound; small positive x gives positive k.\\
372--380&PASS&Physical realization preserves effective coefficients and the equilibrium equation.\\
381--388&PASS&Large G enforces a positive and avoids at most one frequency zero; no universal existence claim survives.\\
389--399&PASS&P07 constants and frequency are conditional five-node inputs only.\\
400--412&PASS&Exact tuple and critical state are strictly positive and physically valid.\\
413--426&PASS&All five labeled eigenvalues, individual simplicity, ordering and exact gap are correct.\\
427--434&PASS&Second root lies strictly between zero and one.\\
435--442&PASS&Exact positivity and original-model residual are checked before spectral reduction.\\
443--452&PASS&Effective data and raw diagonal formula agree independently.\\
453--465&PASS&Radical real parts and strict noncritical difference are correct; table corroboration is redundant.\\
466--474&PASS&Frequency signs and exact positive bounds yield the complete complex spectrum.\\
475--484&PASS&P02 polynomial factors exactly; no P10 input is used.\\
485--492&PASS&Positive lower root and discriminant establish the upper-root identification.\\
493--504&PASS&Frozen g-path and unique local root near one are accurately scoped.\\
505--514&PASS&Explicit vector uses positive exponents, normalized pairing and correct adjoint sign.\\
515--526&PASS&Real block and both spectral-projector inverse sums are correct.\\
527--534&PASS&Mode-zero and mode-two denominators have exact values and are nonzero.\\
535--545&PASS&At X=1 the bracket derivative equals lambda\_0; local positivity and root ordering persist.\\
546--555&PASS&Vector and denominator proofs introduce no tensors or nonlinear conventions.\\
556--563&PASS&H1-H2 are proved physical and original-field assertions.\\
564--565&PASS&H3 is local smooth upper-branch existence only.\\
566--569&PASS&H4-H6 assert the exact pair, individual simplicity and stable gap.\\
570--571&PASS&H7-H8 are exactly the two proved spectral exclusions.\\
572--573&PASS&H9 records established vectors, pairing and plane orientation.\\
574--575&PASS&H10 explicitly leaves transversality and path derivatives unestablished.\\
576--577&PASS&H11 explicitly leaves Lyapunov and nonlinear nondegeneracy unestablished.\\
578--580&PASS&H12 withholds periodic orbits, branch stability and rotating waves.\\
581--587&PASS&Ledger references correctly distinguish results from exclusions.\\
588--592&PASS&Exact nonemptiness is analytical; computation is corroboration and C073 is unused.\\
593--596&PASS&B02/B06 spectral choices are scoped; B07 nonlinear normalization stays with P13.\\
597--600&PASS&First loss/open nonlinear regime remains unproved; proposed note is separate and frozen files untouched.\\
\bottomrule\end{longtable}

\section{Request coverage and artifact inventory}
The source-level ledger above is supplemented by this task-check mapping.
The numbers refer to the supplied audit request, not P11 equation numbers.
\begin{longtable}{p{0.17\linewidth}p{0.75\linewidth}}
\toprule Request items&Disposition and location\\\midrule\endhead
0--1&Current-version authentication and actual dependency use: Section~\ref{sec:auth} and the hash inventory.\\
2&All six layer verdicts are separately stated, with an overall verdict.\\
3--7&Orientation, reachability, frequency, strict definition and full label coverage: Section~\ref{sec:A}.\\
8--14&Simplicity, gap, real dimensions, inverse gates, adjoint, real block, ties and repeated zero: Section~\ref{sec:B}.\\
15--19&Equilibrium chart, real-part function, k graph, physical transfer and rank-two locus: Section~\ref{sec:C}.\\
20--29&Independent necessity/sufficiency, physical realization, frequency avoidance and exact obstructions: Section~\ref{sec:D}.\\
30--46&Original-field residual, exact spectrum, upper root, local branch, vectors, block and denominators: Section~\ref{sec:E}.\\
47--50&Literal H1--H12 ledger, separate clarification recommendation and P12/P13 boundaries: Section~\ref{sec:F}.\\
51--52&Analytical counterexample attempts and independent exact/finite checks: Section~\ref{sec:checks}.\\
53&Complete 1--600 range partition: Section~\ref{sec:lines} and companion CSV.\\
54--58&Required verdicts, exact slug, audit-only ZIP, source/PDF, preservation and explicit closure handoff.\\
\bottomrule
\end{longtable}
The audit-only archive is named exactly
\path{cyclic_network_project_A11_audit_only.zip}, with sole project root
\path{cyclic_network_project/}. It contains this standalone audit source,
its PDF reading copy and \path{audits/A11_support/} evidence only.
It contains no replacement P11 source and is not a full-project snapshot.
The support files record the audited version; all requested input hashes;
all 200 original project hashes; archive identity and CRC checks; independent
checker and results; the 646-cell realizability table; complete line review;
clarification review; validation provenance; and preservation/build records.
The original draft archives and every pre-existing project file are retained.

\section{Authenticated input hashes}
The full machine-readable inventory includes byte counts, line counts and
proof-body hashes. The following human-readable list records all 31 inputs.
Each displayed SHA-256 is split only for page width; concatenate its halves.

\par\medskip\noindent\begin{minipage}{\linewidth}
\noindent\path{project_manifest.md}\par
\noindent{\small\ttfamily e57b131d77a821547cc30a2b1efb21e1\\4a71548d6d40668a50140d9bc8e17272}\par
\end{minipage}\par

\par\medskip\noindent\begin{minipage}{\linewidth}
\noindent\path{claim_registry.md}\par
\noindent{\small\ttfamily 4e1a7fd2a38b4364e98570c55b211c27\\7348082693ab1e56367222e4f4d37b8f}\par
\end{minipage}\par

\par\medskip\noindent\begin{minipage}{\linewidth}
\noindent\path{frozen_theorem_package.md}\par
\noindent{\small\ttfamily 724adf3d04a7f9661a2e53fbae359192\\e8a5a5ee56737bdca2379209b556ea63}\par
\end{minipage}\par

\par\medskip\noindent\begin{minipage}{\linewidth}
\noindent\path{final_novelty_verdict.md}\par
\noindent{\small\ttfamily 5fa8ddd5ada9628b104dba68dd535356\\442c5e62ec45b72f3c2c354dec6037ea}\par
\end{minipage}\par

\par\medskip\noindent\begin{minipage}{\linewidth}
\noindent\path{novelty_decisions.md}\par
\noindent{\small\ttfamily 8fd079fb108b93c4504bc2b527de11c1\\f7e867a967c7cef0d90e1ff7a0a9175e}\par
\end{minipage}\par

\par\medskip\noindent\begin{minipage}{\linewidth}
\noindent\path{open_decisions.md}\par
\noindent{\small\ttfamily 506c634c9958d13b4d6ef4080ebe214a\\f5fd0c77d5dc8fe6f1b91d2750ffab59}\par
\end{minipage}\par

\par\medskip\noindent\begin{minipage}{\linewidth}
\noindent\path{02_theorem_dependency_map.tex}\par
\noindent{\small\ttfamily a1a4435f36fc1d1b6f0b2c5226a91328\\3c76658d55b96493e6c50a8693831cac}\par
\end{minipage}\par

\par\medskip\noindent\begin{minipage}{\linewidth}
\noindent\path{proofs_compendium.tex}\par
\noindent{\small\ttfamily 88afbd023e84e6396f34b785fe54dffe\\4fcc398796e6ff3ec6da45300776b62f}\par
\end{minipage}\par

\par\medskip\noindent\begin{minipage}{\linewidth}
\noindent\path{shared_macros.tex}\par
\noindent{\small\ttfamily f4a18a440e1214b619bf18f4a81fee50\\7239460ca796d8ca3916791069d6a7ad}\par
\end{minipage}\par

\par\medskip\noindent\begin{minipage}{\linewidth}
\noindent\path{00_canonical_model.tex}\par
\noindent{\small\ttfamily c209981cb2924e8a782f24aa2904e8c8\\5a105ff99319dd6aada12f48a351c2a4}\par
\end{minipage}\par

\par\medskip\noindent\begin{minipage}{\linewidth}
\noindent\path{01_notation_and_assumptions.tex}\par
\noindent{\small\ttfamily 97d993c6d3986524dbcc29659ee73d25\\7be1b03a35c44401e8ef4534303eeb91}\par
\end{minipage}\par

\par\medskip\noindent\begin{minipage}{\linewidth}
\noindent\path{proofs/P01_wellposedness_and_absorbing_box.tex}\par
\noindent{\small\ttfamily 0d6950821aa83f3d09e8bfe84dab2b0b\\5020aa6e7881cf6928e984bff1890f1e}\par
\end{minipage}\par

\par\medskip\noindent\begin{minipage}{\linewidth}
\noindent\path{audits/A01_wellposedness_and_absorbing_box.tex}\par
\noindent{\small\ttfamily 8d5c1b2ec782dbe49b07d2bbe529ee47\\4d4c142695c53b523ccb41d63bb33098}\par
\end{minipage}\par

\par\medskip\noindent\begin{minipage}{\linewidth}
\noindent\path{proofs/P02_synchronous_subspace_and_equilibria.tex}\par
\noindent{\small\ttfamily 8e057c2cd009659eee45906844e25f4f\\38c156e5eb2e4130e824aa5e8299c51a}\par
\end{minipage}\par

\par\medskip\noindent\begin{minipage}{\linewidth}
\noindent\path{audits/A02_synchronous_subspace_and_equilibria.tex}\par
\noindent{\small\ttfamily ca3d723d9fcfac14f5129583441e611a\\65e41848ec44f3b5e3a0072f795d84ac}\par
\end{minipage}\par

\par\medskip\noindent\begin{minipage}{\linewidth}
\noindent\path{proofs/P03_origin_and_linearization.tex}\par
\noindent{\small\ttfamily 6c3c9bfb2307bdc043eb7f6a874cb222\\e2ac2ce0750b339181f0088dd0bd859e}\par
\end{minipage}\par

\par\medskip\noindent\begin{minipage}{\linewidth}
\noindent\path{audits/A03_origin_and_linearization.tex}\par
\noindent{\small\ttfamily e1008f98ba9667b92b8d74f1bc7d21e0\\21414b454ac976ae0db1500b4b7ef8d1}\par
\end{minipage}\par

\par\medskip\noindent\begin{minipage}{\linewidth}
\noindent\path{proofs/P04_arbitrary_nq_fourier_spectrum.tex}\par
\noindent{\small\ttfamily 1331c8598acd0481046250d918cac028\\ca6d13d2604612448c3459e864321c58}\par
\end{minipage}\par

\par\medskip\noindent\begin{minipage}{\linewidth}
\noindent\path{audits/A04_arbitrary_nq_fourier_spectrum.tex}\par
\noindent{\small\ttfamily 8a57079241269c92a3c3e8dc4dccf859\\2da2d1a636fca0c0990a4f318396d869}\par
\end{minipage}\par

\par\medskip\noindent\begin{minipage}{\linewidth}
\noindent\path{proofs/P05_q2_mode_selection.tex}\par
\noindent{\small\ttfamily 054a047a55ef37815485f6c225be2f63\\22ca8cf1c272614dd1dc4bb10ae8e1dd}\par
\end{minipage}\par

\par\medskip\noindent\begin{minipage}{\linewidth}
\noindent\path{audits/A05_q2_mode_selection.tex}\par
\noindent{\small\ttfamily 0f7e9836b91009a201185a4acbbb9971\\5f3c65df715e405962b7d6852ff27031}\par
\end{minipage}\par

\par\medskip\noindent\begin{minipage}{\linewidth}
\noindent\path{proofs/P07_five_node_specialization.tex}\par
\noindent{\small\ttfamily 1258641da2b1b729523da6262cf75345\\f5ef1eabf7a9f693a8b13dd34deec393}\par
\end{minipage}\par

\par\medskip\noindent\begin{minipage}{\linewidth}
\noindent\path{audits/A07_five_node_specialization.tex}\par
\noindent{\small\ttfamily b3923d6520cc249c508b80dd06a2b515\\f054e1510e255224a864e986cd38ce4b}\par
\end{minipage}\par

\par\medskip\noindent\begin{minipage}{\linewidth}
\noindent\path{proofs/P08_matched_bilinear_analysis.tex}\par
\noindent{\small\ttfamily 81d273f0bad1714d2470e9fc826f0819\\e9dde367a9c42633705a275b404e5454}\par
\end{minipage}\par

\par\medskip\noindent\begin{minipage}{\linewidth}
\noindent\path{audits/A08_matched_bilinear_analysis.tex}\par
\noindent{\small\ttfamily 87cd4a03061648eb4cf1c65d1afa47c7\\132c1f7eed10d546784ee166c8e5c935}\par
\end{minipage}\par

\par\medskip\noindent\begin{minipage}{\linewidth}
\noindent\path{proofs/P09_transcritical_bifurcation.tex}\par
\noindent{\small\ttfamily 5d357a5be97dfbad62a0107ca2e010ff\\97e0c11a719ee850df69411275f11ff8}\par
\end{minipage}\par

\par\medskip\noindent\begin{minipage}{\linewidth}
\noindent\path{audits/A09_transcritical_bifurcation.tex}\par
\noindent{\small\ttfamily 1234e02d7cf620511bfd7b9306dbb72a\\c68929ef01e8c3ba70bd3615bcd2b943}\par
\end{minipage}\par

\par\medskip\noindent\begin{minipage}{\linewidth}
\noindent\path{proofs/P10_saddle_node_and_bistability.tex}\par
\noindent{\small\ttfamily 37542cc5409d7db33ca183869933fbdb\\ad673628f69f6988f2778a84bc84094f}\par
\end{minipage}\par

\par\medskip\noindent\begin{minipage}{\linewidth}
\noindent\path{audits/A10_saddle_node_and_bistability.tex}\par
\noindent{\small\ttfamily ff2ebb0f84ae79fc9aad012a3945ef23\\6e8ccdac9f47d64fd708e4dc66963b19}\par
\end{minipage}\par

\par\medskip\noindent\begin{minipage}{\linewidth}
\noindent\path{proofs/P11_hopf_spectral_conditions.tex}\par
\noindent{\small\ttfamily 4be8e9b1e64c022b635f3cca9f7f49ff\\bc5358a2250dda1715cb60e82d9b073b}\par
\end{minipage}\par

\par\medskip\noindent\begin{minipage}{\linewidth}
\noindent\path{PROPOSED_P11_scope_clarification.md}\par
\noindent{\small\ttfamily 027c130a0ed4e472335746689f63a194\\d286fd01648558b04feb18d1f3132b0a}\par
\end{minipage}\par

\section{Closure handoff}
\textbf{OVERALL VERDICT: VERIFIED. Required corrections: none.}
There are no E01, E02 or other required mathematical or wording issues to
close. This verdict applies only to the exact P11 v1.0 hash identified at
the start of this audit. The mathematical body requires no change.

The \textbf{ORIGINAL P11 development chat may perform administrative closure
without changing the mathematical body}. It should preserve this exact
version identity, record the independent A11 verdicts in the appropriate
lifecycle records, retain the scope exclusions and validation-provenance
distinctions, and handle the separately recommended planning clarification
explicitly. This audit does not itself mark P11 VERIFIED and does not edit
P11, its manifest, the registry, the compendium, or a frozen file.

After successful correction/closure, the original P11 development chat
must create \path{cyclic_network_project_P11_audited.zip} as the
\textbf{COMPLETE full-project snapshot}. That archive is not created in
this audit chat. The existing \path{cyclic_network_project_P11_draft.zip}
must remain unchanged. The deliverable here is solely the audit-only archive.

\textbf{P12 is cleared to begin after administrative closure.} This is
permission to undertake its scoped proof work, not certification of a P12
result. P12 must analyze the actual frozen physical path and account for
equilibrium motion; this audit computes no derivative along it. P13 must
still establish nonlinear nondegeneracy with reconciled conventions before
any assembled Hopf, periodic-orbit or branch-stability conclusion. First-loss
and rotating-wave claims remain subject to their separate outstanding gates.
Neither P12 nor P13 is begun here.
\end{document}
